\section{Lezione 10} \hfill 17/03/2026
\subsection{Castagneto per legname da opera (segati)}
Intuitivamente, per ottenere assortimenti legnosi da opera, servirebbe un governo a fustaia e non a ceduo. Purtroppo, dopo il primo taglio di una fustaia si ottiene un popolamento a ceduo.\\
La soluzione è gestire un ceduo come se fosse una fustaia. Al fine di ottenere diametri superiori ai 30 \unit{cm} (diametro minimo per l'utilizzo in segheria), è possibile perseguire 2 strade:
\begin{itemize}
    \item selvicoltura di popolamento: l'obiettivo è avere 300-500 fusti ad ha.\\
    Le migliaia di fusti che si hanno all'inizio del turno sono normalmente concentrati sulle ceppaie: occorre lavorare sulle ceppaie con una serie di diradamenti, che portino all'obiettivo finale del numero di piante.\\
    La scelta dei candidati avviene nell'ordine:
    \begin{enumerate}
        \item le eventuali piante nate da seme (se c'è posto solitamente ci sono sempre);
        \item le ceppaie dominanti;
        \item andare a conclusione con le ceppaie codominanti (se non ci sono abbastanza fusti tra i precedenti).
    \end{enumerate}
    In passato, al 2\unit{^\circ}-3\unit{^\circ} anno, veniva fatto uno sfollo sulle ceppaie, dimezzando il numero di polloni e mantenendo i migliori. Oggi non viene più fatto a causa del cancro corticale, che attacca le ferite della ceppaia, e che ammazzerebbe i polloni.\\
    Perciò, il primo diradamento viene svolto dopo i 6-10 anni, anche in modo che le piante da seme raggiungano le grandezze dei polloni.\\
    È necessario non tagliare più del 50 percento dei polloni presenti in una ceppaia, poichè sennò la ceppaia ricaccerebbe in maniera eccessiva.\\
    Ogni 4-5 anni si riduce il numero di polloni, fino ad ottenere il numero finale (300-500 ad ettaro), in modo che ogni ceppaia abbia 1-2 polloni.\\
    Tutte le operazioni vanno terminate entro i 20 anni di età, perché il castagno (come tutte le fagacee ad esclusione del faggio), in condizione di stress tende ad emettere rami epicornici (falsi polloni sulla chioma), che riducono la qualità del legname. Prima dei 20 anni, la ceppaia risponde allo stress emettendo polloni, dopo i 20 anni emettendo rami epicornici.\\
    Dopo i 20 anni da inizio del turno, si aspetta il raggiungimento del diametro obiettivo, e poi si effettua il taglio.
    \item selvicoltura d'albero: intorno ai 5-7 anni dall'inizio del turno, vengono eletti i candidati (al massimo 100-150 ad ettaro, di elevatissima qualità), scelti tra:
    \begin{enumerate}
        \item le piante nate da seme;
        \item i polloni.
    \end{enumerate} 
    Questi candidati vengono liberati man mano dalla competizione; il resto del popolamento svolge la funzione di protezione dei candidati e/o per i classici assortimenti prodotti dal castagno. 
\end{itemize}
Il cancro corticale oggi si presenta in due forme: la virulenta che si diffonde subito e la ipovirulenta che si diffonde dopo.\\
Il castagno è come ``il maiale", si usa tutto, infatti:
\begin{itemize}
    \item le radici hanno una importante funzione di produzione dei funghi porcini;
    \item la corteccia è utilizzata per produrre tannino;
    \item dal tronco si producono assortimenti legnosi;
    \item le foglie venivano usate come lettiera per gli animali;
    \item i frutti venivano/sono usati per il consumo alimentare.
\end{itemize}
Il castagno è stato portato in diversi luoghi/aree, anche oltre i suoi limiti, grazie alla sua versatilità.\\
L'estensione attuale dei castagneti in Italia equivale a 700000 ha:
\begin{itemize}
    \item tali popolamenti non dovrebbero naturalmente esistere (essendo il castagno una specie sporadica);
    \item le stazioni sono diverse: variano tra le poco alle molto fertili.
\end{itemize}
In passato c'è stata una fase di abbandono di parte del castagno, causa l'abbandono della campagna a favore della città (soprattutto per la fascia di colline difficili da coltivare e poco produttive).\\
Cosa fare di questi castagneti abbandonati? 
\begin{itemize}
    \item Se non conviene economicamente/storicamente/strutturalmente riprenderli, è meglio lasciarli andare.\\
    Il destino di questi popolamenti è quello di sparire; il castagno invece rimane una specie sporadica. Il popolamento viene sostituito dalle specie che dovrebbero essere presenti (es. siano queste querce, faggio, latifoglie miste).\\
    In questo caso dunque non serve fare niente.\\
    Le stazioni pessime per il castagno vengono abbandonate.
   \item Nelle stazioni fertili, e dove c'è interesse per il castagno, il bosco degradato viene recuperato.\\
    Nei castagneti da legno, basta tagliare il castagno, in modo che possa ripartire con i propri ricacci.\\
    Dato il mercato ancora attivo delle castagne, i castagneti da frutto sono stati inizialmente ripresi per la produzione del frutto, ed ora sono gestiti anche per questioni turistico-ricreative.\\
\end{itemize}
Il castagneto da frutto abbandonato è normalmente attaccato dal cancro; normalmente ci sono 20-30 piante ad ettaro ed all'interno c'è crescita di altre specie. Il castagneto da frutto è riconoscibile dalla presenza di piante innestate.\\
La ripresa del popolamento si articola in diverse fasi:
\begin{itemize}
    \item eliminare tutto ciò che non è castagno, lasciando le piante del vecchio popolamento da frutto ed eventualmente il selvaggione di castagno cresciuto dentro;
    \item è così possibile capire la condizione di salute del popolamento (``le piante sono recuperabili o sono morte?'');
    \item se non ci sono sufficienti piante, occorre piantarne di nuove: è possibile usare il selvaggione come porta innesto (sono già presenti ed eradicati) e/o piantare nuovi individui;
    \item dagli anni successivi occorre fare una pulizia ogni anno.
\end{itemize}
L'operazione di ripresa è costosa, ma permette una nuova fruttificazione dopo tre-quattro anni.\\ 
Il castagno è più aggressivo della robinia (che quest'ultima ha una chioma meno leggera).\\
\subsection{Bagolaro - Celtis australis}
È una formazione particolare, non tanto presente in Veneto, ma molto in Piemonte, Lombardia e Friuli Venezia Giulia (dove il substrato è acido).\\
È questa specie (anche detta ``spaccasassi'') è competitiva nelle aree dei roccioni, delle aree ripide, dove cresce in tasche di suolo povere. Se cresce in un buon substrato può diventare alto fino a 30 \unit{m}.\\
Ha un buon legno sia per legna da ardere che per la costruzione dei carri (per il fatto che è molto elastico).\\
Veniva gestito con ceduazioni di turno pari a 20-30 anni.
Non è una specie interessante per la selvicoltura boschiva; bensì è molto utilizzata in selvicoltura urbana, perché resiste bene all'inquinamento.
\subsection{Betulla - Betula}
Questo popolamento si pone nel passaggio dal piano collinare verso quello superiore (e ci entra anche dentro).\\
In Italia ci sono 3 tipologie di betulla: Betula alba (bianca o verrucosa), Betula pubescens (anche detta pelosa, si trova a Nord delle Alpi e nelle zone più umide, si differenzia con quella italiana solamente per i rametti grigio-glabri), e Betula etnensis (dell'Etna).\\
Questa specie:
\begin{itemize}
    \item è lo stereotipo della specie pioniera (sia primaria che secondaria);
    \item è fondamentalmente acidofila (Lombardia e Piemonte ne hanno molta anche in pianura);
    \item è estremamente opportunista;
    \item l'apparato radicale è profondo;
    \item ha una vita breve (50-70 anni); 
    \item ha un seme leggerissimo ed una fioritura molto precoce (intorno a febbraio-marzo);
    \item i tempi di fioritura, fruttificazione e dispersione del seme molto lunghi;
    \item i semi, prodotti in grandissima quantità, sono trilobati;
    \item la disseminazione è anemocora.
\end{itemize}
È di interesse selvicolturale perchè ha occupato gran parte dei pascoli abbandonati, essendo molto opportunista: la sua capacità caulinare non è grandiosa, ma ha una grande disseminazione da seme.\\
Sulle Alpi ci sono almeno 3000 ha di betuleti, che possono essere utilizzati in vari modi. La betulla è sempre presente nei boschi misti di latifoglie o nei confini dei boschi di castagno.\\
Il legno di betulla:
\begin{multicols}{2}
\begin{itemize}
    \item è molto bianco;
    \item si lavora bene ed è molto stabile (es. compensati marini, cucine);
    \item è molto usato per i mobili e per la produzione di stuzzicadenti;
    \item è di media-bassa densità;
    \item è uno dei migliori legnami da ardere, poiché ricchissimo di oli essenziali;
    \item ha il difetto di non essere durabile (marcisce nel giro di una stagione, non è adatto per usi esterni). Se è legna da ardere dev'essere ben protetta.
\end{itemize}
\end{multicols}
I popolamenti di betulla possono essere usati come:
\begin{itemize}
    \item nursering per i boschi definitivi (es. querceta, faggeta);
    \item riserve di biomassa o legnami da opera: mediante gestione a taglio raso o taglio raso con riserve, in modo che ogni 30-40 anni possano essere ottenute grandi quantità di legno di buona qualità oppure legname da sfoglia/sega.
\end{itemize}
\subsection{Ontani bianchi}
Tra collinare e montana\\
Specie azonale, ontanete di ontano bianco.\\
Tra i quattro ontani presenti in Italia, ontano nero, napoletano, bianco e di pianura....1:03\\
Gli ontani sono legati all'acqua. Essendo in bassa montagna, gli ontaneti di ontano bianco indicano che c'è acqua, eccessiva per le altre specie. Si trovano lungo i torrenti, vicino alle sorgenti, (di fondovalle, di canalone o di versante), 1:05. Queste formazioni sono importanti dal punto di vista della biodiversità. Quelli di canalone non sono molto utilizzati, ma sono un indicatore del tempo di ritorno della valanga: ontaneti giovani indicano il recente passaggio di una valanga.\\
Le ontanete di versante sono quasi scomparse causa la captazione delle fontane create: quelle rimaste devono essere mantenute.\\
Quelle di fondovalle sono quasi scomparse, causa la conversione spesso al meleto.\\
Ha un legno molto morbido\\
In passato, gli ontaneti erano importanti per la vita di montagna: erano la fonte di legna da ardere (di buona qualità) più vicina a casa; erano i primi posti dove potevano essere portati i capi di bestiame in estate. \\
Queste formazioni sono quasi del tutto scomparse.\\
La gestione da svolgere è il ceduo semplice; non bisogna che le piante invecchino troppo
\subsection{Pioppo tremulo - Populus tremula}
I fusti sono molto dritti, cresce molto in fretta.\\
La chioma è leggera e le foglie sono tremolanti, con picciolo lungo.\\
Terreno umido.\\
capacità pollonifera caulinare e fortemente radicale\\
I popolamenti sono molto densi nelle aree umide. Non si lascia moltiplicare per talee (unico pioppo). Non viene coltivato\\
Vita breve (come tutti i pioppi).\\
\subsection{Acero-frassineti e acero tiglieti}
In contesto 400-800-1000 metri di quota.\\
Sono composti dall'acero di montano, acero riccio, acero oparo (appenninico), acero campestre (più di pianura). I tigli sono cordata e platifillo. Il cordato ha le foglie più piccole e la parte gialla alla base del picciolo. Il tiglio usato in città sono normalmente degli ibridi.\\
Queste specie sono le cosiddette latifoglie nobili, e hanno un comportamento ecologico molto simile tra di loro. Sono specie normalmente meso-file (riesco a rinnovare in qualsiasi condizione). SOno pioniere secondarie (necessitano di suoli ricchi e già formati) = pioniere secondarie della faggeta. hanno accrescimento abbastanza veloce (pioniere), i legni sono duri. SOno le classiche specie da suoli freschi, sciolti e profondi: rifornimento costante d'acqua, facile da penetrare, necessità di approfondire gli apparati radicali.\\
Gli acero-frassineti e gli acero-  non sono formazioni misti, ma sono dominate a momenti alterni tra le due specie. \\
In natura gli acero frassineti sono ... entra la definitiva (normalmente il faggio) ed  oppure sono popolazioni stabili (es forre canyon), dove c'è forte umidità atmosferica e sacche costanti\\
Il frassino è usato come (legno duro, molto elastico, molto chiaro appena tagliato, si ingallisce con l'ossidazione) in passato usato per le cucime . era usato per i macini degli attrezzi agricoli (picconi, pale). I primi sci erano in frassino; oggi i migliori sci hanno l'anima in legno di frassino. venivano diffusi in filari o piante isolate in gran parte dei maggenghi (alpeggi di maggio): in questi maggenghi poteva nevicare, l'erba spariva dal suolo, il frassino aveva già fogliato ed i rami venivano sgamollati per dare in pasto al bestiame. La frasca di frassino venivano anche usate per migliorare la buonificazione della panna (più facile fare il burro).  